\chapter{Etude du chatbot Django + ML}

\section{Objectif du chatbot}

Le projet chatbot fournit un service de question/reponse local pour assister les utilisateurs sur des sujets academie (inscription, frais, planning, contact, etc.).

\section{Stack technique}

Le fichier \texttt{requirements.txt} indique une base Django 5.2.1 completee par des outils data science:
\begin{itemize}[leftmargin=*]
  \item \texttt{pandas}
  \item \texttt{numpy}
  \item \texttt{scikit-learn}
  \item \texttt{rapidfuzz}
\end{itemize}

\section{Exposition API}

Le routage principal est defini par:
\begin{itemize}[leftmargin=*]
  \item \texttt{/} pour l'interface web;
  \item \texttt{/chat} pour chat web;
  \item \texttt{/api/chat} pour integration applicative externe.
\end{itemize}

Le contrat \texttt{/api/chat} retourne notamment:
\begin{itemize}[leftmargin=*]
  \item \texttt{response}
  \item \texttt{score}
  \item \texttt{category}
  \item \texttt{source}
  \item \texttt{matched\_question}
\end{itemize}

\section{Moteur de reponse}

Le service \texttt{ml\_index.py} implemente une strategie en trois niveaux:
\begin{enumerate}[leftmargin=*]
  \item exact match sur question normalisee;
  \item fuzzy match (RapidFuzz) si exact non trouve;
  \item recherche TF-IDF + cosine similarity via scikit-learn.
\end{enumerate}

Le service \texttt{bot.py} combine aussi des intents predefinis avant de passer sur la recherche ML.

\section{Configuration et tuning}

Le fichier \texttt{settings.py} expose:
\begin{itemize}[leftmargin=*]
  \item \texttt{QA\_CSV}: chemin du dataset question/reponse;
  \item \texttt{MIN\_SIM}: seuil minimal similarity;
  \item \texttt{FUZZY\_MIN}: seuil fuzzy matching;
  \item \texttt{CHATBOT\_API\_KEY}: protection optionnelle de \texttt{/api/chat}.
\end{itemize}

\section{Securite API}

Le module \texttt{api\_auth.py} permet de forcer une cle API via header \texttt{X-API-Key}. Cela facilite l'isolation de l'endpoint en environnement partage.

\section{Integration avec Flutter}

Le frontend Flutter utilise \texttt{ChatbotService.dart}:
\begin{itemize}[leftmargin=*]
  \item envoi \texttt{POST} JSON vers \texttt{CHATBOT\_API\_CHAT\_URL};
  \item gestion du timeout via \texttt{API\_TIMEOUT};
  \item support optionnel de \texttt{CHATBOT\_API\_KEY}.
\end{itemize}

Cette integration garde le chatbot decouple du backend principal, ce qui simplifie son evolution independante.
